\documentclass[10pt,a4paper]{amsart}
\usepackage[margin=19mm]{geometry}
\usepackage{amsmath,amssymb,mathtools}
\usepackage[scr=rsfs,cal=euler]{mathalfa}
\usepackage{xfrac}
\usepackage[unicode,hidelinks,bookmarks=false]{hyperref}
\newcommand{\ie}{i.e.\ }
\newcommand{\cf}{cf.\ }
\newcommand{\resp}{respectively\ }
\newcommand{\Subsection}{\subsection}
\providecommand{\nobreakdash}{\nobreak\hbox{-}}
\setlength{\emergencystretch}{2em}
\multlinegap0pt
\usepackage{bm}
\usepackage{bbm}
\usepackage{dsfont}
\usepackage{xspace}
\usepackage{cancel}
\usepackage{mathtools}
\usepackage{amsthm}
\makeatletter
\@ifundefined{lem}{\newtheorem{lem}{Lemma}}{}
\makeatother
\newcommand{\Imagine}{{\rm Im \hspace{0.05cm} }}
\newcommand{\Real}{{\rm Re \hspace{0.05cm} }}
\newcommand{\dive}{{\rm div \hspace{0.05cm} }}
\title[Removable singularity of (-1)-homogeneous solutions of stati]{Removable singularity of (-1)-homogeneous solutions of stationary Navier-Stokes equations}
\author{Li Li, YanYan Li, Xukai Yan}
\date{Source: arXiv:2410.11170v2 (CC BY 4.0)}
\begin{document}
\maketitle

\noindent\emph{Source and licence.} Removable singularity of (-1)-homogeneous solutions of stationary Navier-Stokes equations. Version arXiv:2410.11170v2 (CC BY 4.0), distributed under the Creative Commons Attribution 4.0 International licence, as recorded in the arXiv licence metadata for this exact version and shown on the abstract page https://arxiv.org/abs/2410.11170v2. Full rights notice: licenses/arxiv-2410.11170-CC-BY-4.0.txt.\par\smallskip

\noindent\textit{Editorial scope.} Counted unit: the Lem whose source label is \texttt{lem:2\_2} in the article (source0.tex lines 248--265, source label \texttt{lem:2\_2}), delivered together with the article's own complete original proof of it (source0.tex lines 266--449, one \texttt{proof} environment of 184 lines). No mathematics is written, repaired or re-derived: statement and proof are the article's own text line for line. Required context retained uncounted: none. Every definition, display and label of those lines is retained unchanged. Omitted from this package and not redistributed: 1--247, 450--1441. Nothing inside a retained excerpt is omitted or altered: every retained body is delivered exactly as the article's lines are written, so no omission receipt of any kind accompanies this package. Citations: the retained excerpts carry no citation command at all, so no bibliography entry of the article is transcribed and no reference is redistributed. Reference adaptation: none was needed and none was made; every reference inside the delivered unit resolves inside it, so no unresolved reference or citation is produced. Printed slips kept literal: no printed word, symbol, hypothesis or number of the counted statement or of its proof was altered. Layout and class: the article's own class is \texttt{article} with packages including fullpage, amsmath, amssymb, amsthm, graphicx, placeins, extpfeil, enumerate, caption, subcaption, hyperref, xcolor, pbox, natbib; the wrapper is the lane's pinned \texttt{amsart} editorial wrapper at 10pt on A4, which restates the theorem-like environments the retained text needs and adds the article's own macro definitions that the retained text needs (\texttt{\textbackslash Imagine}, \texttt{\textbackslash Real}, \texttt{\textbackslash dive}), with extra wrapper packages (bm, bbm, dsfont, xspace, cancel, mathtools). The delivered numbering is the unit's own. Assets: no retained span contains a figure environment, a graphics-inclusion command, a PDF-page inclusion, a table or a TikZ picture; no image file is copied into this package, the package declares no asset dependency and no image file is redistributed with it. Licence: version arXiv:2410.11170v2 of this article is distributed under the Creative Commons Attribution 4.0 International licence, as recorded in the arXiv licence metadata for this exact version and shown on the abstract page https://arxiv.org/abs/2410.11170v2. A full rights notice is delivered at licenses/arxiv-2410.11170-CC-BY-4.0.txt. Characters: every retained excerpt is pure ASCII, so the delivered engine, XeLaTeX with its default Latin Modern text font, reports no missing character.
\begin{lem}\label{lem:2_2}
	Let $R>0$, 
	$a<b$ satisfying $ab>0$,
	and $F\in C^{\infty}(\bar{\Omega}_{a, b, R}\setminus\{x'=0\})$ be a $(-3)$-homogeneous vector-valued function.
	Suppose $q\in C^{\infty}(\bar{\Omega}_{a, b, R}\setminus\{x'=0\})$ is a $(-2)$-homogeneous function satisfying 
	\begin{equation*}
		-\Delta q=\dive F(x), 
		\quad \textrm{ in }\Omega_{a, b, R}\setminus\{x'=0\}. 
	\end{equation*}
	Assume there exists some $\delta\in (0, 2)$ such that 
	\begin{equation}\label{eq:lem:2_2:1}
		|q(x)| |x'|^{2-\delta}+\sum_{j=0}^{2}|\nabla^jF| |x'|^{j+2-\delta}=o(1)
	\end{equation}
	as $x'\to 0$ uniformly in $\Omega_{a, b, R}\setminus\{x'=0\}$. Then there exist some $h(x) \in L^{\infty}_{loc}((a, b), W^{1, s}(D_R))$ for any $1< s<\frac{2}{2-\delta}$, and $a_0(x_3), a_1(x_3), b_1(x_3)\in C(a, b)$, such that 
	\begin{equation}\label{eq:lem:2_2:2}
		q(x)=h(x)+a_0(x_3)\ln |x'|+a_1(x_3)\frac{x_1}{|x'|^2}+b_1(x_3)\frac{x_2}{|x'|^2}, \quad \textrm{ in }\Omega_{a, b, R}\setminus\{x'=0\}.
	\end{equation}
\end{lem}

\begin{proof}
	We prove the lemma when $a<b<0$. The proof when $0<a<b$ is similar. 
	
	Let $a', b'$ be arbitrary numbers satisfying $a<a'<b'<b$. For any fixed $\bar{x}\in \Omega_{a', b', R/2} \setminus\{x'=0\}$, let $\bar{\rho}:=\min\{|\bar{x}'|/3, (R-|\bar{x}'|)/3, (b-b')/3, (a'-a)/3\}$, so $B_{\bar{\rho}}(\bar{x})\subset \Omega_{a, b, R}$. Set $\bar{q}(y)=\bar{\rho}^{2-\delta}q(\bar{x}+\bar{\rho}y)$ and $\bar{F}(y)=\bar{\rho}^{3-\delta}F(\bar{x}+\bar{\rho}y)$ for $y\in B_2$. By (\ref{eq:lem:2_2:1}), we have $\sup_{B_2}(|\bar{q}|+\sum_{j=0}^{2}|\nabla^j\bar{F}|)=o(1)$ as $\bar{\rho}\to 0$. Moreover, 
	\[
		-\Delta \bar{q}(y)=\dive \bar{F},
		\quad y\in B_2. 
	\] 
	By elliptic theories, $\sup_{B_1}(|\nabla \bar{q}|+|\nabla^2 \bar{q}|)=o(1)$ as $\bar{\rho}\to 0$. Therefore, 
	\begin{equation}\label{eq:lem:2_2:pf:1}
		\sum_{j=1}^{2}|\nabla^j q(x)| |x'|^{j+2-\delta}=o(1), 
		\textrm{ as }x'\to 0, \textrm{ uniformly for }x_3\in (a', b'). 
	\end{equation}
	Denote $\Delta' = \partial_1^2 + \partial_2^2$. For each fixed $x_3\in (a, b)$, we have
	\[
		\Delta' q(x',x_3)=-\dive F(x', x_3)-\partial_3^2q(x',x_3)=:g_0(x', x_3)+g_1(x', x_3), 
	\]
	where $g_0(x', x_3):=-\partial_1F_1(x', x_3)-\partial_2F_2(x', x_3)$ 
	and $g_1(x', x_3):=-\partial_3F_3(x', x_3)-\partial_3^2q(x',x_3)$. 
	
	We first study the existence and regularity of the solutions $q_0$ and $q_1$ of the Poisson equation 
	\begin{equation}\label{eq:Poisson}
	\left\{
	\begin{aligned}
		& \Delta' q_i(x', x_3)=g_i(x', x_3), \quad x'\in D_R, \\
		& q_i(\cdot, x_3)|_{\partial D_R}=0,
	\end{aligned}
	\right.
	\end{equation}
	for $i=0,1$, and then estimate the remaining part $q_2:=q-q_0-q_1$. 
	
	(1) Since $F=o(|x'|^{\delta-2})$, we have 
	 \[
	 \sup_{a<x_3<b}\|g_0(\cdot, x_3)\|_{W^{-1, s}(D_R)}=\sup_{a<x_3<b} \| F(\cdot, x_3) \|_{L^{s}(D_{R})}<\infty, \quad \forall 1< s<\frac{2}{2-\delta}. 
	 \]
	So for each $x_3\in (a, b)$, there exists a solution $q_0(\cdot, x_3)\in W^{1, s}(D_R)$ of (\ref{eq:Poisson}) for $i=0$. Moreover, we have $\sup_{a<x_3<b} \|q_0(\cdot, x_3) \|_{W^{1, s}(D_{R})}<\infty$.

	Since $F$ is $(-3)$-homogeneous, $F_3(x)=|x_3|^{-3}F_3(-x_1/x_3, -x_2/x_3, -1)$. Without loss of generality, assume $a'<-1<b'$. 
	Then by (\ref{eq:lem:2_2:1}), we have
	\[
		|\partial_3F_3(x', x_3)|
		\le C(|F_3(-\frac{x}{x_3})|+ |x'| |\nabla F_3(-\frac{x}{x_3})|)=o\left(\frac{1}{|x'|^{2-\delta}}\right) 
	\]
	as $x'\to 0$, uniformly for $x_3\in (a, b)$, 
	where $C$ is some constant depending only on $a, b, R$. 
Moreover, since $q$ is $(-2)$-homogeneous, we have $q(x)=|x_3|^{-2}q(-x_1/x_3, -x_2/x_3, -1)$ in $\Omega_{a, b, R}\setminus\{x'=0\}$. 
	Then by (\ref{eq:lem:2_2:1}) and (\ref{eq:lem:2_2:pf:1}), we have
	\[
		|\partial_3^2q(x', x_3)|
		\le C(|q(-\frac{x}{x_3})|+|x'| |\nabla q(-\frac{x}{x_3})|+|x'|^2|\nabla^2 q(-\frac{x}{x_3})|)=o\left(\frac{1}{|x'|^{2-\delta}}\right)
			\]
	as $x'\to 0$, uniformly for $x_3\in (a', b')$, where $C$ is some constant depending only on $a, b, a', b', R$. 
	Thus we have $g_1(x', x_3)=o\left(|x'|^{\delta-2}\right)$ uniformly for $x_3\in (a', b')$ and 	 $\sup_{a'<x_3<b'}\|g_1(\cdot, x_3)\|_{L^s(D_R)}<\infty$ for any $1< s<\frac{2}{2-\delta}$. 
	In particular, since $a', b'$ are arbitrary numbers in $(a,b)$, we have $\|g_1(\cdot, x_3)\|_{L^s(D_R)}<\infty$ for any $x_3\in (a, b)$. Thus for each $x_3\in (a, b)$, there exists a solution $q_1(\cdot,x_3)\in W^{2,s}(D_R)$ of (\ref{eq:Poisson}) for $i=1$. Moreover, we have $\sup_{a'<x_3<b'} \| q_1(\cdot, x_3) \|_{W^{2,s}(D_{R})} < \infty$ for any $a<a'<b'<b$ and $1< s<\frac{2}{2-\delta}$, since $\sup_{a'<x_3<b'}\|g_1(\cdot, x_3)\|_{L^s(D_R)}<\infty$. So $q_0, q_1$ are well-defined in $\Omega_{a, b, R}$ and 
	\begin{equation}\label{eq:lem:2_2:pf:2}
		q_0, q_1 \in L^{\infty}_{loc}((a, b), W^{1, s}(D_R)), \quad \forall 1<s<\frac{2}{2-\delta}.
	\end{equation} 
	Since $g_0, g_1\in C^{\infty}(\bar{\Omega}_{a, b, R}\setminus\{x'=0\})$, we also have 
	$q_0, q_1\in C^{\infty}(\bar{\Omega}_{a, b, R}\setminus\{x'=0\})$. 	

	(2) Let $q_2(x',x_3):=q(x',x_3)-(q_0(x', x_3)+q_1(x', x_3))$, then $\Delta' q_2(x', x_3)=0$ in $\Omega_{a, b, R}\setminus\{x'=0\}$, and for any $a<a'<b'<b$, 
	\begin{equation}\label{eq:lem:2_2:pf:3}
		\sup\limits_{\substack{R/4<|x'|<R \\ a'<x_3<b'}}(|q_2|+|\nabla q_2|)<\infty. 
	\end{equation}
	We will show 
	\begin{equation}\label{eq:lem:2_2:pf:4}
		q_2(x)=\hat{h}(x)+a_0(x_3)\ln |x'|+a_1(x_3)\frac{x_1}{|x'|^2}+b_1(x_3)\frac{x_2}{|x'|^2}, \quad x\in \Omega_{a, b, R}\setminus\{x'=0\},
	\end{equation}
	for some $\hat{h}(x) \in L^{\infty}_{loc}((a, b), W^{1, s}(D_R))$ for any $1< s<\frac{2}{2-\delta}$ and $a_0(x_3), a_1(x_3), b_1(x_3)\in C(a, b)$. Then (\ref{eq:lem:2_2:2}) follows from (\ref{eq:lem:2_2:pf:2}) and (\ref{eq:lem:2_2:pf:4}) by setting $h=\hat{h}+q_0+q_1$. 
	
	For any $f(x')\in C^{\infty}(D_R\setminus\{0\})$, denote 
	\begin{equation}\label{eq:lem:2_2:pf:5}
		\hat{f}(x'):=f(x')-\left(\frac{1}{2\pi}\int_{\partial D_{R/2}}\frac{\partial f}{\partial \nu}\right)\ln |x'|.
	\end{equation}
	For each $x_3\in (a, b)$, let $\hat{q}_2(x', x_3)$ be defined as above. Since $\Delta' q_2=0$ in $D_R\setminus\{0\}$, we have $\Delta' \hat{q}_2=0$ in $D_R\setminus\{0\}$ and $\int_{\partial D_{R'}}\frac{\partial \hat{q}_2}{\partial \nu}=0$ for any $0<R'<R$ and $a<x_3<b$. Let $z=x_1+ix_2$, and
	\[
		w(z, x_3): 
		=\hat{q}_2+i \int_{(\frac{R}{2}, 0)}^z(-\partial_2\hat{q}_2 dx_1+\partial_1\hat{q}_2dx_2), 
	\]
	where the integral $\int_{(\frac{R}{2}, 0)}^{z}$ is independent of the path in $D_R\setminus\{0\}$. 
	Then $w$ is analytic in $z$ in $\{0<|z|< R\}$. 
	The Laurent series of $w$ in $z$ takes the form $w(z, x_3)=\sum_{m=-\infty}^{\infty} c_m(x_3) z^m$. 
	For any $0<R_1<R_2<R$, the series is uniformly convergent in $\{ R_1\le |z|\le R_2\}$. For 
	any $\rho>0$, we have 
	\begin{equation}\label{eq:lem:2_2:pf:6}
		c_m(x_3)=\frac{1}{2\pi i}\oint_{|z|=\rho} \frac{w(z, x_3)}{z^{m+1}}. 
	\end{equation}
	Now we show that $c_m(x_3)\equiv 0$ in $(a', b')$ for all $m\le -2$, 
	and $\sup_{a'<x_3<b'}|c_{-1}(x_3)|<\infty$. 
	
	\noindent\emph{Claim}: 
	$\sup_{a'<x_3<b'} \|w(\cdot, x_3)\|_{L^{s}(D_{R})}<\infty$ for any $1<s<\frac{2}{2-\delta}$.
	
	We will prove the claim later. Suppose the claim holds. Note that (\ref{eq:lem:2_2:pf:6}) holds for any $\rho>0$, we have
	\[
		c_{m}(x_3)=\frac{1}{2\pi i}\oint_{|z|=\rho} \frac{w(z, x_3)}{z^{m+1}}=\frac{1}{\pi i \rho}\int_{\rho/2}^{\rho}\oint_{|z|=t} \frac{w(z, x_3)}{z^{m+1}}dS(z)dt=\frac{1}{\pi i \rho}\int_{D_{\rho}\setminus D_{\rho/2}}\frac{w(z, x_3)}{z^{m+1}}.
	\]
	Now fix some $1<s<\frac{2}{2-\delta}$. By the Claim and H\"{o}lder's inequality, we have
	\[
		\sup_{a'<x_3<b'}|c_{m}(x_3)|\le \frac{1}{\pi\rho}\sup_{a'<x_3<b'}\|w\|_{L^s(D_{\rho}\setminus D_{\rho/2})}\|z^{-m-1}\|_{L^{\frac{s}{s-1}}(D_{\rho}\setminus D_{\rho/2})} 
		\le C\rho^{-m-\frac{2}{s}}
	\]
	for all $\rho>0$ and some $C$ independent of $\rho$. 
	Since $s>1$, we have $-m-2/s>-m-2$. By sending $\rho\to 0$, we have $\sup_{a'<x_3<b'}|c_m(x_3)|=0$ for $m\le -2$. As $a<a', b'<b$ are arbitrary, we have $c_m(x_3)\equiv 0$ for $x_3\in (a, b)$ and $m\le -2$. Then $w(z, x_3)=\sum_{m=-1}^{\infty}c_m(x_3)z^m$. By the definition of $w$ and (\ref{eq:lem:2_2:pf:3}), we have $\sup_{\substack{R/4<|x'|<R \\ a'<x_3<b'}}|w|<\infty$. Taking $\rho=R/2$ in (\ref{eq:lem:2_2:pf:6}), we have $\sup_{a'<x_3<b'}|c_m(x_3)|\le C(R/2)^{-m}$ for $m\ge -1$ for some $C$ depending only on $\sup_{\substack{R/4<|x'|<R \\ a'<x_3<b'}}|w|$. So $\sup_{a'<x_3<b'}|c_{-1}(x_3)|<\infty$. Moreover, for $|z|\le R/4$, 
	\begin{equation}\label{eq:lem:2_2:pf:7}
		| \sum_{m=0}^{\infty} c_m(x_3) z^m | \le C \sum_{m=0}^{\infty} \frac{1}{2^m}<\infty. 
	\end{equation}
	Set $\hat{h}(x', x_3)=\Real (\sum_{m=0}^{\infty}c_m(x_3)z^m)$, $a_0(x_3)=\frac{1}{2\pi}\int_{\partial D_{R/2}}\frac{\partial q_2}{\partial \nu}$, $a_1(x_3)=\Real c_{-1}(x_3)$ and $b_1(x_3)=\Imagine c_{-1}(x_3)$. By the fact that 
	\[
		\Real w(z, x_3)=\hat{q}_2=q_2(x', x_3)-(\frac{1}{2\pi}\int_{\partial D_{R/2}}\frac{\partial q_2}{\partial \nu})\ln |x'|, 
	\]
	we have for any $x_3\in(a,b)$ that
	\[
	\begin{split}
		q_2(x', x_3) & =\Real w(z, x_3)+(\frac{1}{2\pi}\int_{\partial D_{R/2}}\frac{\partial q_2}{\partial \nu})\ln |x'|\\
		& =\Real (\sum_{m=-1}^{\infty}c_m(x_3)z^m)+(\frac{1}{2\pi}\int_{\partial D_{R/2}}\frac{\partial q_2}{\partial \nu})\ln |x'|\\
		& =\hat{h}(x', x_3)+
		a_0(x_3)\ln |x'|+a_1(x_3)\frac{x_1}{|x'|^2}+b_1(x_3)\frac{x_2}{|x'|^2}. 
		\end{split}
	\]
	So (\ref{eq:lem:2_2:pf:4}) is proved. 
	
	Let $h:=\hat{h}+q_0+q_1$, then
	\[
		q=(q_0+q_1)+q_2=h(x', x_3)+a_0(x_3)\ln |x'|+a_1(x_3)\frac{x_1}{|x'|^2}+b_1(x_3)\frac{x_2}{|x'|^2}. 
	\]
	By (\ref{eq:lem:2_2:pf:2}) and (\ref{eq:lem:2_2:pf:7}), we have $h\in L^{\infty}_{loc}((a, b), W^{1, s}(D_R))$. By (\ref{eq:lem:2_2:pf:3}), we have $a_0\in L^{\infty}_{loc}(a, b)$. Since $\sup_{a'<x_3<b'}|c_{-1}(x_3)|\le C$, we have $a_1, b_1\in L^{\infty}_{loc}(a, b)$. Since $q\in C^{\infty}(\bar{\Omega}_{a, b, R}\setminus\{x'=0\})$, we have $q_1, q_2, \hat{q}_2\in C^{\infty} (\Omega_{a, b, R}\setminus\{x'=0\})$. Thus for each $z\ne 0$, $w(z, x_3)$ is continuous in $x_3\in (a, b)$ and $c_m$ is continuous in $x_3\in (a, b)$, and therefore $a_0(x_3), a_1(x_3)$, $b_1(x_3)\in C(a, b)$. 

	(3)	\noindent\emph{Proof of Claim}: 
	Recall that $q_2=q-(q_0+q_1)$, we have $\Real w=\hat{q}_2=\hat{q}-(\hat{q}_1+\hat{q}_0)$, where $\hat{q}(\cdot, x_3), \hat{q}_0(\cdot, x_3), \hat{q}_1(\cdot, x_3)$ are defined by (\ref{eq:lem:2_2:pf:5}) for each $x_3\in (a', b')$. By (\ref{eq:lem:2_2:1}) and (\ref{eq:lem:2_2:pf:2}), we have 
	\[
		\sup_{a'<x_3<b'}\|\hat{q}_2\|_{L^s(D_R)}
		<\infty, \quad \forall 1<s<\frac{2}{2-\delta}. 
	\]
	It remains to show 
	\begin{equation}\label{eq:lem:2_2:pf:8}
		\sup_{a'<x_3<b'} \|\Imagine w\|_{L^{s}(D_{R})}<\infty, \quad \forall 1<s<\frac{2}{2-\delta}.
	\end{equation} 
	Note $\Imagine w=\int_{(\frac{R}{2}, 0)}^z(-\partial_2\hat{q}_2 dx_1+\partial_1\hat{q}_2dx_2)$. For any $z\in D_R\setminus \{z_2=0\}$, denote $\bar{z}=\frac{R}{2}\frac{z}{|z|}$. Let $\Gamma_1$ be the counter-clockwise path from $(R/2, 0)$ to $\bar{z}$ along $\partial D_{R/2}$, and $\Gamma_2$ be the path from $\bar{z}$ to $z$ along the ray in the direction of $z$, and let $\Gamma=\Gamma_1\cup\Gamma_2$. For any $f=(f_1, f_2)\in C(D_R\setminus\{0\})$, define
	\begin{equation*}
		\mathcal{L}[f](z) := \int_{\Gamma}(f_1dx_1+f_2dx_2), \quad z\in D_R\setminus [0, R). 
	\end{equation*}
	Then the following facts hold.

	\noindent\emph{Fact 1}. If $f=(f_1, f_2)\in C(D_R\setminus\{0\})$ satisfies $|f_1(z)|+|f_2(z)|\le C_0|z|^{\lambda}$ for some $C_0>0$ and $\lambda\in\mathbb{R}\setminus\{-1\}$, then $|\mathcal{L}[f](z)|\le C(1+|z|^{\lambda+1})$ for some $C$ depending only on $R, \lambda$ and $C_0$. 
	
	To see this, let $(r, \theta)$ be the polar coordinates in $\mathbb{R}^2$, where $x_1=r\cos\theta$ and $x_2=r\sin\theta$. For any $z\in D_R\setminus \{z_2=0\}$, denote $z=(|z|, \theta_0)$, we have 
	\[
	\begin{split}
		|\mathcal{L}[f](z)| & =|\int_{\Gamma}(f_1\sin\theta+f_2\cos\theta)rd\theta+(f_1\cos\theta+f_2\sin\theta)dr|\\
		& \le \frac{R}{2}\int_{0}^{\theta_0}|f\big(\frac{R}{2}, \theta\big)|d\theta 
		+\int_{|z|}^{\frac{R}{2}}|f(r, \theta_0)|dr\\
		& 
		\le C+C|\int_{|z|}^{\frac{R}{2}}r^{\lambda}dr| 
		\le C+C|z|^{\lambda+1}.
	\end{split}
	\]
	
	\noindent\emph{Fact 2.} If $f=(f_1, f_2)\in C(D_R\setminus\{0\})\cap L^s(D_R)$ for some $s\ge 1$, then $\mathcal{L}[f]\in L^s(D_R)$, and 
	\[
		\|\mathcal{L}[f]\|_{L^s(D_R)}\le C(1+\|f\|_{L^{s}(D_R)}) 
	\]
	for some $C$ depending only on $s$ and $R$. 
	
	As in the proof of Fact 1, we have 
	\[
	\begin{split}
		|\mathcal{L}[f](z)| & \le \frac{R}{2}\int_{0}^{\theta_0}|f\big(\frac{R}{2}, \theta\big)|d\theta 
		+\int_{|z|}^{\frac{R}{2}}|f(r, \theta_0)|dr
		\le C+C\big|\int_{|z|}^{\frac{R}{2}}|f(r, \theta_0)|^sdr\big|^{\frac{1}{s}}. 
	\end{split}
	\]
	Taking the power $s$ of the above and integrating in $z$ over $D_R$, we have 
	\[
		\|\mathcal{L}[f](z)\|_{L^s(D_R)}\le C(1+\|f\|_{L^{s}(D_R)}). 
	\]
	So Fact 2 holds. 
	
	Since in the definition of $w$, the integral $\int_{(\frac{R}{2}, 0)}^{z}$ is independent of path, we take the path to be $\Gamma$ as defined above. Then 
	\[
		\Imagine w=\int_{\Gamma}(-\partial_2\hat{q}_2 dx_1+\partial_1\hat{q}_2dx_2)=\mathcal{L}[\nabla^{\perp}\hat{q}_2]=\mathcal{L}[\nabla^{\perp}\hat{q}]-\mathcal{L}[\nabla^{\perp}(\hat{q}_0+\hat{q}_1)].
	\]
	where $\nabla^{\perp}=(-\partial_2, \partial_1)$. By (\ref{eq:lem:2_2:pf:1}) and Fact 1, we have $\sup_{a'<x_3<b'}|z|^{2-\delta}|\mathcal{L}[\nabla^{\perp}\hat{q}(\cdot, x_3)]|<\infty$, and therefore $\sup_{a'<x_3<b'} \|\mathcal{L}[\nabla^{\perp}\hat{q}(\cdot, x_3)]\|_{L^{s}(D_{R})}<\infty$ for $1<s<\frac{2}{2-\delta}$. By (\ref{eq:lem:2_2:pf:2}) and Fact 2, we have $\sup_{a'<x_3<b'} \|\mathcal{L}[\nabla^{\perp}\hat{q}_0(\cdot, x_3)+\nabla^{\perp}\hat{q}_1(\cdot, x_3)]\|_{L^{s}(D_{R})}<\infty$ for $1<s<\frac{2}{2-\delta}$. Thus (\ref{eq:lem:2_2:pf:8}) holds and the Claim is proved. The lemma is proved. 
\end{proof}

\end{document}
